%%%%%%%%%%%%%%%%%%%%%%%%%%%%%%%%%%%%%%%%%%%%%%%%%%%%%%%%%%%%
% 8.17 SURVIVAL ANALYSIS
% The Composition of Advanced Mathematics
%%%%%%%%%%%%%%%%%%%%%%%%%%%%%%%%%%%%%%%%%%%%%%%%%%%%%%%%%%%%

\section{Survival Analysis}
\label{sec:survival-analysis}

\prerequisite{
Probability, random variables, probability distributions, expectation,
conditional probability, statistical estimation, confidence intervals,
hypothesis testing, regression, generalized linear models, and basic
calculus.
}

\learningobjectives{
By the end of this section, the reader should be able to:
\begin{itemize}
    \item explain why time-to-event data require specialized methods;
    \item define survival time and event time;
    \item distinguish complete, censored, and truncated observations;
    \item distinguish right, left, and interval censoring;
    \item define the survival function;
    \item define the cumulative distribution function for event times;
    \item define the density function for continuous survival times;
    \item define the hazard function;
    \item define the cumulative hazard function;
    \item derive relationships among survival, hazard, density, and
          cumulative hazard functions;
    \item interpret hazard rates correctly;
    \item understand common parametric survival distributions;
    \item work with exponential survival models;
    \item understand Weibull survival models;
    \item estimate survival curves using the Kaplan--Meier estimator;
    \item understand risk sets;
    \item understand censored contributions to survival estimation;
    \item estimate cumulative hazards using the Nelson--Aalen estimator;
    \item compare survival curves using the log-rank test;
    \item understand proportional hazards;
    \item understand the Cox proportional-hazards model;
    \item interpret Cox regression coefficients as hazard ratios;
    \item understand the partial likelihood concept;
    \item recognize time-dependent covariates;
    \item recognize nonproportional hazards;
    \item understand stratified Cox models conceptually;
    \item recognize accelerated failure-time models;
    \item distinguish hazard-ratio and time-ratio interpretations;
    \item understand competing risks conceptually;
    \item distinguish cause-specific hazards and cumulative incidence;
    \item understand recurrent-event data conceptually;
    \item recognize frailty models;
    \item understand informative and noninformative censoring;
    \item understand survival prediction and median survival;
    \item recognize discrete-time survival analysis;
    \item connect survival analysis with reliability, regression,
          stochastic processes, and statistical learning.
\end{itemize}
}

%%%%%%%%%%%%%%%%%%%%%%%%%%%%%%%%%%%%%%%%%%%%%%%%%%%%%%%%%%%%
\subsection{What Is Survival Analysis?}
\label{subsec:what-is-survival-analysis}
%%%%%%%%%%%%%%%%%%%%%%%%%%%%%%%%%%%%%%%%%%%%%%%%%%%%%%%%%%%%

Survival analysis studies the time until an event occurs.

The event may be:

\[
\boxed{
\text{death},
}
\]

\[
\boxed{
\text{disease recurrence},
}
\]

\[
\boxed{
\text{machine failure},
}
\]

\[
\boxed{
\text{customer churn},
}
\]

or

\[
\boxed{
\text{completion of a process}.
}
\]

The methods are designed to handle incomplete observation of event times.

%%%%%%%%%%%%%%%%%%%%%%%%%%%%%%%%%%%%%%%%%%%%%%%%%%%%%%%%%%%%
\subsection{Survival Time}
\label{subsec:survival-time}
%%%%%%%%%%%%%%%%%%%%%%%%%%%%%%%%%%%%%%%%%%%%%%%%%%%%%%%%%%%%

Let

\[
\boxed{
T\geq0
}
\]

denote the random time until the event of interest.

The distribution of \(T\) may be described using:

\[
\boxed{
F(t),
}
\]

\[
\boxed{
S(t),
}
\]

\[
\boxed{
f(t),
}
\]

and

\[
\boxed{
h(t).
}
\]

%%%%%%%%%%%%%%%%%%%%%%%%%%%%%%%%%%%%%%%%%%%%%%%%%%%%%%%%%%%%
\subsection{Event-Time Distribution Function}
\label{subsec:event-time-cdf}
%%%%%%%%%%%%%%%%%%%%%%%%%%%%%%%%%%%%%%%%%%%%%%%%%%%%%%%%%%%%

The cumulative distribution function of \(T\) is

\[
\boxed{
F(t)
=
P(T\leq t).
}
\]

It gives the probability that the event has occurred by time \(t\).

Since

\[
T\geq0,
\]

typically

\[
F(0)
\]

is zero or small depending on the model.

%%%%%%%%%%%%%%%%%%%%%%%%%%%%%%%%%%%%%%%%%%%%%%%%%%%%%%%%%%%%
\subsection{Survival Function}
\label{subsec:survival-function}
%%%%%%%%%%%%%%%%%%%%%%%%%%%%%%%%%%%%%%%%%%%%%%%%%%%%%%%%%%%%

\begin{definitionbox}{Survival Function}
The survival function is

\[
\boxed{
S(t)
=
P(T>t).
}
\]
\end{definitionbox}

Since

\[
F(t)
=
P(T\leq t),
\]

we have

\[
\boxed{
S(t)
=
1-F(t).
}
\]

%%%%%%%%%%%%%%%%%%%%%%%%%%%%%%%%%%%%%%%%%%%%%%%%%%%%%%%%%%%%
\subsection{Basic Properties of the Survival Function}
\label{subsec:survival-function-properties}
%%%%%%%%%%%%%%%%%%%%%%%%%%%%%%%%%%%%%%%%%%%%%%%%%%%%%%%%%%%%

The survival function satisfies

\[
\boxed{
0\leq S(t)\leq1.
}
\]

It is nonincreasing.

For many continuous survival models,

\[
\boxed{
S(0)=1,
}
\]

and

\[
\boxed{
\lim_{t\to\infty}
S(t)
=
0.
}
\]

%%%%%%%%%%%%%%%%%%%%%%%%%%%%%%%%%%%%%%%%%%%%%%%%%%%%%%%%%%%%
\subsection{Density of a Continuous Survival Time}
\label{subsec:survival-density}
%%%%%%%%%%%%%%%%%%%%%%%%%%%%%%%%%%%%%%%%%%%%%%%%%%%%%%%%%%%%

If \(T\) is continuous with density \(f\), then

\[
\boxed{
f(t)
=
F'(t).
}
\]

Since

\[
S(t)
=
1-F(t),
\]

we obtain

\[
\boxed{
f(t)
=
-
S'(t).
}
\]

Thus the survival function completely determines the density.

%%%%%%%%%%%%%%%%%%%%%%%%%%%%%%%%%%%%%%%%%%%%%%%%%%%%%%%%%%%%
\subsection{Hazard Function}
\label{subsec:hazard-function}
%%%%%%%%%%%%%%%%%%%%%%%%%%%%%%%%%%%%%%%%%%%%%%%%%%%%%%%%%%%%

\begin{definitionbox}{Hazard Function}
The hazard function is

\[
\boxed{
h(t)
=
\lim_{\Delta t\to0^+}
\frac{
P
(
t\leq T<t+\Delta t
\mid
T\geq t
)
}{
\Delta t
}.
}
\]
\end{definitionbox}

For a continuous survival time,

\[
\boxed{
h(t)
=
\frac{
f(t)
}{
S(t)
}.
}
\]

%%%%%%%%%%%%%%%%%%%%%%%%%%%%%%%%%%%%%%%%%%%%%%%%%%%%%%%%%%%%
\subsection{Interpretation of the Hazard}
\label{subsec:hazard-interpretation}
%%%%%%%%%%%%%%%%%%%%%%%%%%%%%%%%%%%%%%%%%%%%%%%%%%%%%%%%%%%%

The hazard describes the instantaneous event rate at time \(t\),
conditional on survival up to that time.

It is not generally a probability.

A hazard may exceed \(1\) because it is a rate per unit time rather than a
probability over a fixed interval.

Thus:

\[
\boxed{
h(t)
\neq
P(T=t).
}
\]

%%%%%%%%%%%%%%%%%%%%%%%%%%%%%%%%%%%%%%%%%%%%%%%%%%%%%%%%%%%%
\subsection{Cumulative Hazard Function}
\label{subsec:cumulative-hazard}
%%%%%%%%%%%%%%%%%%%%%%%%%%%%%%%%%%%%%%%%%%%%%%%%%%%%%%%%%%%%

The cumulative hazard is

\[
\boxed{
H(t)
=
\int_0^t
h(u)\,du.
}
\]

Because

\[
h(t)
=
-
\frac{
S'(t)
}{
S(t)
},
\]

integration gives

\[
\boxed{
H(t)
=
-
\log S(t).
}
\]

Therefore,

\[
\boxed{
S(t)
=
e^{-H(t)}.
}
\]

%%%%%%%%%%%%%%%%%%%%%%%%%%%%%%%%%%%%%%%%%%%%%%%%%%%%%%%%%%%%
\subsection{Density from Hazard}
\label{subsec:density-from-hazard}
%%%%%%%%%%%%%%%%%%%%%%%%%%%%%%%%%%%%%%%%%%%%%%%%%%%%%%%%%%%%

Since

\[
f(t)
=
h(t)S(t),
\]

and

\[
S(t)
=
e^{-H(t)},
\]

we obtain

\[
\boxed{
f(t)
=
h(t)
e^{-H(t)}.
}
\]

Thus the hazard function also determines the full survival distribution.

%%%%%%%%%%%%%%%%%%%%%%%%%%%%%%%%%%%%%%%%%%%%%%%%%%%%%%%%%%%%
\subsection{Core Survival Relationships}
\label{subsec:core-survival-relationships}
%%%%%%%%%%%%%%%%%%%%%%%%%%%%%%%%%%%%%%%%%%%%%%%%%%%%%%%%%%%%

For continuous survival times,

\[
\boxed{
S(t)=1-F(t),
}
\]

\[
\boxed{
f(t)=-S'(t),
}
\]

\[
\boxed{
h(t)=\frac{f(t)}{S(t)},
}
\]

\[
\boxed{
H(t)=\int_0^t h(u)\,du,
}
\]

and

\[
\boxed{
S(t)=e^{-H(t)}.
}
\]

These identities form the mathematical core of survival analysis.

%%%%%%%%%%%%%%%%%%%%%%%%%%%%%%%%%%%%%%%%%%%%%%%%%%%%%%%%%%%%
\subsection{Worked Example: Survival from a Hazard}
\label{subsec:worked-survival-from-hazard}
%%%%%%%%%%%%%%%%%%%%%%%%%%%%%%%%%%%%%%%%%%%%%%%%%%%%%%%%%%%%

\begin{workedexamplebox}{Constant Hazard}

Suppose

\[
h(t)=0.2
\]

for all \(t\geq0\).

Then

\[
H(t)
=
\int_0^t0.2\,du
=
0.2t.
\]

Therefore,

\[
S(t)
=
e^{-0.2t}.
\]

\begin{answerbox}
\[
\boxed{
S(t)
=
e^{-0.2t}.
}
\]
\end{answerbox}

\end{workedexamplebox}

%%%%%%%%%%%%%%%%%%%%%%%%%%%%%%%%%%%%%%%%%%%%%%%%%%%%%%%%%%%%
\subsection{Expected Survival Time}
\label{subsec:expected-survival-time}
%%%%%%%%%%%%%%%%%%%%%%%%%%%%%%%%%%%%%%%%%%%%%%%%%%%%%%%%%%%%

For a nonnegative continuous random variable,

\[
\boxed{
\mathbb{E}[T]
=
\int_0^\infty
S(t)\,dt,
}
\]

provided the integral is finite.

This identity is often convenient because survival functions are central
to event-time models.

%%%%%%%%%%%%%%%%%%%%%%%%%%%%%%%%%%%%%%%%%%%%%%%%%%%%%%%%%%%%
\subsection{Median Survival Time}
\label{subsec:median-survival-time}
%%%%%%%%%%%%%%%%%%%%%%%%%%%%%%%%%%%%%%%%%%%%%%%%%%%%%%%%%%%%

The median survival time \(m\) satisfies

\[
\boxed{
S(m)=0.5.
}
\]

Equivalently,

\[
\boxed{
F(m)=0.5.
}
\]

Median survival is often more robust and easier to estimate than mean
survival when the upper tail is incompletely observed.

%%%%%%%%%%%%%%%%%%%%%%%%%%%%%%%%%%%%%%%%%%%%%%%%%%%%%%%%%%%%
\subsection{Why Ordinary Methods Fail}
\label{subsec:why-ordinary-survival-methods-fail}
%%%%%%%%%%%%%%%%%%%%%%%%%%%%%%%%%%%%%%%%%%%%%%%%%%%%%%%%%%%%

Suppose a study ends before every subject experiences the event.

For some individuals, the true event time is unknown.

If these observations are discarded, information is wasted and estimates
can be biased.

If their last observed follow-up times are treated as actual event times,
the event-time distribution is distorted.

Survival methods explicitly incorporate this incomplete information.

%%%%%%%%%%%%%%%%%%%%%%%%%%%%%%%%%%%%%%%%%%%%%%%%%%%%%%%%%%%%
\subsection{Censoring}
\label{subsec:censoring}
%%%%%%%%%%%%%%%%%%%%%%%%%%%%%%%%%%%%%%%%%%%%%%%%%%%%%%%%%%%%

Censoring occurs when the exact event time is not observed, but partial
information about it is known.

The most common form is right censoring.

%%%%%%%%%%%%%%%%%%%%%%%%%%%%%%%%%%%%%%%%%%%%%%%%%%%%%%%%%%%%
\subsection{Right Censoring}
\label{subsec:right-censoring}
%%%%%%%%%%%%%%%%%%%%%%%%%%%%%%%%%%%%%%%%%%%%%%%%%%%%%%%%%%%%

A survival time is right censored if one knows only that

\[
\boxed{
T>C.
}
\]

For example, a subject may still be event-free at the end of the study.

Then the censoring time \(C\) provides a lower bound on the unknown event
time.

%%%%%%%%%%%%%%%%%%%%%%%%%%%%%%%%%%%%%%%%%%%%%%%%%%%%%%%%%%%%
\subsection{Left Censoring}
\label{subsec:left-censoring}
%%%%%%%%%%%%%%%%%%%%%%%%%%%%%%%%%%%%%%%%%%%%%%%%%%%%%%%%%%%%

A survival time is left censored if one knows only that

\[
\boxed{
T<C.
}
\]

For example, an event may already have occurred before the first
observation.

%%%%%%%%%%%%%%%%%%%%%%%%%%%%%%%%%%%%%%%%%%%%%%%%%%%%%%%%%%%%
\subsection{Interval Censoring}
\label{subsec:interval-censoring}
%%%%%%%%%%%%%%%%%%%%%%%%%%%%%%%%%%%%%%%%%%%%%%%%%%%%%%%%%%%%

A survival time is interval censored if it is only known that

\[
\boxed{
L<T\leq R.
}
\]

This occurs when subjects are checked periodically and the event occurs
between visits.

%%%%%%%%%%%%%%%%%%%%%%%%%%%%%%%%%%%%%%%%%%%%%%%%%%%%%%%%%%%%
\subsection{Observed Right-Censored Data}
\label{subsec:observed-right-censored-data}
%%%%%%%%%%%%%%%%%%%%%%%%%%%%%%%%%%%%%%%%%%%%%%%%%%%%%%%%%%%%

Let

\[
T_i
=
\text{event time},
\]

and

\[
C_i
=
\text{censoring time}.
\]

The observed time is

\[
\boxed{
Y_i
=
\min(T_i,C_i).
}
\]

Define the event indicator

\[
\boxed{
\delta_i
=
\mathbf{1}_{\{T_i\leq C_i\}}.
}
\]

Then

\[
\delta_i=1
\]

means the event was observed, while

\[
\delta_i=0
\]

means the observation was right censored.

%%%%%%%%%%%%%%%%%%%%%%%%%%%%%%%%%%%%%%%%%%%%%%%%%%%%%%%%%%%%
\subsection{Likelihood Contribution Under Right Censoring}
\label{subsec:right-censoring-likelihood}
%%%%%%%%%%%%%%%%%%%%%%%%%%%%%%%%%%%%%%%%%%%%%%%%%%%%%%%%%%%%

For an observed event at time \(y_i\), the likelihood contribution is

\[
f(y_i).
\]

For a censored observation at time \(y_i\), the contribution is

\[
S(y_i),
\]

because the only information is

\[
T_i>y_i.
\]

Thus the combined likelihood is

\[
\boxed{
L
=
\prod_{i=1}^{n}
[
f(y_i)
]^{\delta_i}
[
S(y_i)
]^{1-\delta_i}.
}
\]

%%%%%%%%%%%%%%%%%%%%%%%%%%%%%%%%%%%%%%%%%%%%%%%%%%%%%%%%%%%%
\subsection{Worked Example: Censored Likelihood Contribution}
\label{subsec:worked-censored-likelihood}
%%%%%%%%%%%%%%%%%%%%%%%%%%%%%%%%%%%%%%%%%%%%%%%%%%%%%%%%%%%%

\begin{workedexamplebox}{Observed Event Versus Censoring}

Suppose two subjects are observed.

Subject \(1\) experiences the event at time \(5\).

Subject \(2\) is censored at time \(8\).

Then their likelihood contribution is

\[
f(5)S(8).
\]

\begin{answerbox}
\[
\boxed{
L
\propto
f(5)S(8).
}
\]
\end{answerbox}

\end{workedexamplebox}

%%%%%%%%%%%%%%%%%%%%%%%%%%%%%%%%%%%%%%%%%%%%%%%%%%%%%%%%%%%%
\subsection{Noninformative Censoring}
\label{subsec:noninformative-censoring}
%%%%%%%%%%%%%%%%%%%%%%%%%%%%%%%%%%%%%%%%%%%%%%%%%%%%%%%%%%%%

Many standard survival procedures assume censoring is noninformative.

Informally, after accounting for modeled information, the censoring
mechanism should not reveal additional information about the future event
time.

A common sufficient formulation is conditional independence between

\[
T
\]

and

\[
C
\]

given relevant covariates.

%%%%%%%%%%%%%%%%%%%%%%%%%%%%%%%%%%%%%%%%%%%%%%%%%%%%%%%%%%%%
\subsection{Informative Censoring}
\label{subsec:informative-censoring}
%%%%%%%%%%%%%%%%%%%%%%%%%%%%%%%%%%%%%%%%%%%%%%%%%%%%%%%%%%%%

Censoring is informative when the censoring mechanism is related to event
risk even after accounting for the analysis model.

For example, subjects at high risk may systematically leave a study.

Then ordinary Kaplan--Meier or Cox procedures may be biased.

Special modeling, weighting, or sensitivity analysis may be required.

%%%%%%%%%%%%%%%%%%%%%%%%%%%%%%%%%%%%%%%%%%%%%%%%%%%%%%%%%%%%
\subsection{Truncation}
\label{subsec:survival-truncation}
%%%%%%%%%%%%%%%%%%%%%%%%%%%%%%%%%%%%%%%%%%%%%%%%%%%%%%%%%%%%

Truncation differs from censoring.

With censoring, an individual is included in the dataset but the event
time is only partly observed.

With truncation, some individuals are not observed at all unless their
event times satisfy a sampling condition.

%%%%%%%%%%%%%%%%%%%%%%%%%%%%%%%%%%%%%%%%%%%%%%%%%%%%%%%%%%%%
\subsection{Left Truncation}
\label{subsec:left-truncation}
%%%%%%%%%%%%%%%%%%%%%%%%%%%%%%%%%%%%%%%%%%%%%%%%%%%%%%%%%%%%

Under left truncation, a subject enters the risk set only after surviving
to an entry time.

If entry occurs at time \(L_i\), the subject is observed only if

\[
\boxed{
T_i>L_i.
}
\]

This is also called delayed entry.

Risk sets must reflect this delayed eligibility.

%%%%%%%%%%%%%%%%%%%%%%%%%%%%%%%%%%%%%%%%%%%%%%%%%%%%%%%%%%%%
\subsection{Risk Set}
\label{subsec:risk-set}
%%%%%%%%%%%%%%%%%%%%%%%%%%%%%%%%%%%%%%%%%%%%%%%%%%%%%%%%%%%%

At event time \(t_j\), the risk set contains individuals who are still
under observation and event-free immediately before \(t_j\).

Let

\[
\boxed{
n_j
}
\]

denote the number at risk just before \(t_j\).

Let

\[
\boxed{
d_j
}
\]

denote the number of observed events at \(t_j\).

These quantities are central to nonparametric survival estimation.

%%%%%%%%%%%%%%%%%%%%%%%%%%%%%%%%%%%%%%%%%%%%%%%%%%%%%%%%%%%%
\subsection{Kaplan--Meier Estimator}
\label{subsec:kaplan-meier}
%%%%%%%%%%%%%%%%%%%%%%%%%%%%%%%%%%%%%%%%%%%%%%%%%%%%%%%%%%%%

The Kaplan--Meier estimator of the survival function is

\[
\boxed{
\widehat{S}(t)
=
\prod_{t_j\leq t}
\left(
1-
\frac{
d_j
}{
n_j
}
\right).
}
\]

The product extends over distinct observed event times.

Censoring changes future risk-set sizes but does not create a survival
drop at the censoring time.

%%%%%%%%%%%%%%%%%%%%%%%%%%%%%%%%%%%%%%%%%%%%%%%%%%%%%%%%%%%%
\subsection{Product-Limit Interpretation}
\label{subsec:product-limit}
%%%%%%%%%%%%%%%%%%%%%%%%%%%%%%%%%%%%%%%%%%%%%%%%%%%%%%%%%%%%

At event time \(t_j\),

\[
\boxed{
1-
\frac{
d_j
}{
n_j
}
}
\]

estimates the conditional probability of surviving past \(t_j\), given
survival immediately before \(t_j\).

The overall survival probability is the product of these conditional
survival probabilities.

%%%%%%%%%%%%%%%%%%%%%%%%%%%%%%%%%%%%%%%%%%%%%%%%%%%%%%%%%%%%
\subsection{Worked Example: Kaplan--Meier}
\label{subsec:worked-kaplan-meier}
%%%%%%%%%%%%%%%%%%%%%%%%%%%%%%%%%%%%%%%%%%%%%%%%%%%%%%%%%%%%

\begin{workedexamplebox}{Product-Limit Survival Estimate}

Suppose observed event times produce the following risk-set information:

\[
\begin{array}{c|cc}
t_j & n_j & d_j\\
\hline
2 & 5 & 1\\
4 & 4 & 1\\
7 & 2 & 1
\end{array}
\]

At time \(4\),

\[
\widehat{S}(4)
=
\left(
1-\frac15
\right)
\left(
1-\frac14
\right).
\]

Therefore,

\begin{answerbox}
\[
\boxed{
\widehat{S}(4)
=
\frac45
\cdot
\frac34
=
0.60.
}
\]
\end{answerbox}

\end{workedexamplebox}

%%%%%%%%%%%%%%%%%%%%%%%%%%%%%%%%%%%%%%%%%%%%%%%%%%%%%%%%%%%%
\subsection{Kaplan--Meier Curve}
\label{subsec:kaplan-meier-curve}
%%%%%%%%%%%%%%%%%%%%%%%%%%%%%%%%%%%%%%%%%%%%%%%%%%%%%%%%%%%%

The Kaplan--Meier estimate is a decreasing step function.

It changes only at observed event times.

At censoring times, the curve does not drop, although the number at risk
for later times decreases.

%%%%%%%%%%%%%%%%%%%%%%%%%%%%%%%%%%%%%%%%%%%%%%%%%%%%%%%%%%%%
\subsection{Median Survival from Kaplan--Meier}
\label{subsec:km-median-survival}
%%%%%%%%%%%%%%%%%%%%%%%%%%%%%%%%%%%%%%%%%%%%%%%%%%%%%%%%%%%%

The estimated median survival is the earliest time \(t\) for which

\[
\boxed{
\widehat{S}(t)
\leq
0.5.
}
\]

If the estimated survival curve never falls to \(0.5\), the median
survival is not estimable from the observed follow-up.

%%%%%%%%%%%%%%%%%%%%%%%%%%%%%%%%%%%%%%%%%%%%%%%%%%%%%%%%%%%%
\subsection{Greenwood Variance Formula}
\label{subsec:greenwood-formula}
%%%%%%%%%%%%%%%%%%%%%%%%%%%%%%%%%%%%%%%%%%%%%%%%%%%%%%%%%%%%

A common estimated variance for the Kaplan--Meier estimator is

\[
\boxed{
\widehat{\operatorname{Var}}
[
\widehat{S}(t)
]
=
\widehat{S}(t)^2
\sum_{t_j\leq t}
\frac{
d_j
}{
n_j(n_j-d_j)
}.
}
\]

This is known as Greenwood's formula.

Transformations are often used when constructing survival confidence
intervals so endpoints remain in \([0,1]\).

%%%%%%%%%%%%%%%%%%%%%%%%%%%%%%%%%%%%%%%%%%%%%%%%%%%%%%%%%%%%
\subsection{Cumulative Hazard Estimation}
\label{subsec:cumulative-hazard-estimation}
%%%%%%%%%%%%%%%%%%%%%%%%%%%%%%%%%%%%%%%%%%%%%%%%%%%%%%%%%%%%

Instead of estimating survival directly, one may estimate

\[
H(t).
\]

Because

\[
S(t)
=
e^{-H(t)},
\]

the two approaches are closely connected.

%%%%%%%%%%%%%%%%%%%%%%%%%%%%%%%%%%%%%%%%%%%%%%%%%%%%%%%%%%%%
\subsection{Nelson--Aalen Estimator}
\label{subsec:nelson-aalen}
%%%%%%%%%%%%%%%%%%%%%%%%%%%%%%%%%%%%%%%%%%%%%%%%%%%%%%%%%%%%

The Nelson--Aalen estimator of the cumulative hazard is

\[
\boxed{
\widehat{H}(t)
=
\sum_{t_j\leq t}
\frac{
d_j
}{
n_j
}.
}
\]

Each event time adds an estimated hazard increment

\[
\frac{d_j}{n_j}.
\]

%%%%%%%%%%%%%%%%%%%%%%%%%%%%%%%%%%%%%%%%%%%%%%%%%%%%%%%%%%%%
\subsection{Worked Example: Nelson--Aalen}
\label{subsec:worked-nelson-aalen}
%%%%%%%%%%%%%%%%%%%%%%%%%%%%%%%%%%%%%%%%%%%%%%%%%%%%%%%%%%%%

\begin{workedexamplebox}{Cumulative Hazard}

Suppose the risk-set information is

\[
(n_1,d_1)
=
(10,1),
\]

\[
(n_2,d_2)
=
(8,2).
\]

Then after the second event time,

\[
\widehat{H}
=
\frac1{10}
+
\frac2{8}.
\]

Therefore,

\begin{answerbox}
\[
\boxed{
\widehat{H}
=
0.10+0.25
=
0.35.
}
\]
\end{answerbox}

\end{workedexamplebox}

%%%%%%%%%%%%%%%%%%%%%%%%%%%%%%%%%%%%%%%%%%%%%%%%%%%%%%%%%%%%
\subsection{Kaplan--Meier and Nelson--Aalen Relationship}
\label{subsec:km-na-relationship}
%%%%%%%%%%%%%%%%%%%%%%%%%%%%%%%%%%%%%%%%%%%%%%%%%%%%%%%%%%%%

Kaplan--Meier uses

\[
\prod
\left(
1-\frac{d_j}{n_j}
\right),
\]

while Nelson--Aalen uses

\[
\sum
\frac{d_j}{n_j}.
\]

For small hazard increments,

\[
\log(1-x)
\approx
-x.
\]

Therefore,

\[
\boxed{
\widehat{S}_{KM}(t)
\approx
e^{-\widehat{H}_{NA}(t)}.
}
\]

They are closely related but not exactly identical in finite samples.

%%%%%%%%%%%%%%%%%%%%%%%%%%%%%%%%%%%%%%%%%%%%%%%%%%%%%%%%%%%%
\subsection{Comparing Survival Curves}
\label{subsec:comparing-survival-curves}
%%%%%%%%%%%%%%%%%%%%%%%%%%%%%%%%%%%%%%%%%%%%%%%%%%%%%%%%%%%%

Suppose two or more groups have Kaplan--Meier curves.

Visual separation alone does not provide a formal statistical test.

The log-rank test compares event experience across groups while
accounting for changing risk sets.

%%%%%%%%%%%%%%%%%%%%%%%%%%%%%%%%%%%%%%%%%%%%%%%%%%%%%%%%%%%%
\subsection{Log-Rank Test}
\label{subsec:log-rank-test}
%%%%%%%%%%%%%%%%%%%%%%%%%%%%%%%%%%%%%%%%%%%%%%%%%%%%%%%%%%%%

At event time \(t_j\), suppose:

\[
n_{1j}
=
\text{group 1 at risk},
\]

\[
n_j
=
\text{total at risk},
\]

and

\[
d_j
=
\text{total events}.
\]

Under equal hazards, the expected number of group 1 events is

\[
\boxed{
E_{1j}
=
d_j
\frac{
n_{1j}
}{
n_j
}.
}
\]

%%%%%%%%%%%%%%%%%%%%%%%%%%%%%%%%%%%%%%%%%%%%%%%%%%%%%%%%%%%%
\subsection{Log-Rank Observed Minus Expected}
\label{subsec:log-rank-observed-expected}
%%%%%%%%%%%%%%%%%%%%%%%%%%%%%%%%%%%%%%%%%%%%%%%%%%%%%%%%%%%%

Across event times,

\[
\boxed{
O_1
=
\sum_jd_{1j},
}
\]

and

\[
\boxed{
E_1
=
\sum_jE_{1j}.
}
\]

The log-rank test compares

\[
\boxed{
O_1-E_1
}
\]

with its null variance.

For two groups, a statistic of the form

\[
\boxed{
\frac{
(O_1-E_1)^2
}{
V_1
}
}
\]

is asymptotically

\[
\chi^2_1
\]

under the null.

%%%%%%%%%%%%%%%%%%%%%%%%%%%%%%%%%%%%%%%%%%%%%%%%%%%%%%%%%%%%
\subsection{Interpretation of the Log-Rank Test}
\label{subsec:log-rank-interpretation}
%%%%%%%%%%%%%%%%%%%%%%%%%%%%%%%%%%%%%%%%%%%%%%%%%%%%%%%%%%%%

The log-rank test is particularly sensitive to persistent proportional
differences in hazards.

Its null hypothesis is broadly that the groups have the same survival
experience under the test framework.

When survival curves cross strongly, the log-rank test may have reduced
power or a difficult interpretation.

%%%%%%%%%%%%%%%%%%%%%%%%%%%%%%%%%%%%%%%%%%%%%%%%%%%%%%%%%%%%
\subsection{Weighted Log-Rank Tests Preview}
\label{subsec:weighted-log-rank}
%%%%%%%%%%%%%%%%%%%%%%%%%%%%%%%%%%%%%%%%%%%%%%%%%%%%%%%%%%%%

Weighted versions give different importance to early or late event times.

A generic statistic uses weights

\[
w_j
\]

in sums such as

\[
\boxed{
\sum_j
w_j
(
d_{1j}-E_{1j}
).
}
\]

Different weights target different alternatives.

Weights should ideally be chosen according to a prespecified scientific
question.

%%%%%%%%%%%%%%%%%%%%%%%%%%%%%%%%%%%%%%%%%%%%%%%%%%%%%%%%%%%%
\subsection{Parametric Survival Models}
\label{subsec:parametric-survival-models}
%%%%%%%%%%%%%%%%%%%%%%%%%%%%%%%%%%%%%%%%%%%%%%%%%%%%%%%%%%%%

A parametric survival model assumes a specific family for \(T\).

Examples include:

\[
\boxed{
\text{exponential},
}
\]

\[
\boxed{
\text{Weibull},
}
\]

\[
\boxed{
\text{log-normal},
}
\]

\[
\boxed{
\text{log-logistic},
}
\]

and

\[
\boxed{
\text{gamma-type models}.
}
\]

These models permit direct extrapolation when their assumptions are
reasonable.

%%%%%%%%%%%%%%%%%%%%%%%%%%%%%%%%%%%%%%%%%%%%%%%%%%%%%%%%%%%%
\subsection{Exponential Survival Distribution}
\label{subsec:exponential-survival}
%%%%%%%%%%%%%%%%%%%%%%%%%%%%%%%%%%%%%%%%%%%%%%%%%%%%%%%%%%%%

Suppose

\[
T
\sim
\operatorname{Exponential}(\lambda).
\]

Then

\[
\boxed{
f(t)
=
\lambda e^{-\lambda t},
\qquad
t\geq0.
}
\]

The survival function is

\[
\boxed{
S(t)
=
e^{-\lambda t}.
}
\]

The hazard is

\[
\boxed{
h(t)=\lambda.
}
\]

%%%%%%%%%%%%%%%%%%%%%%%%%%%%%%%%%%%%%%%%%%%%%%%%%%%%%%%%%%%%
\subsection{Memoryless Survival Property}
\label{subsec:memoryless-survival}
%%%%%%%%%%%%%%%%%%%%%%%%%%%%%%%%%%%%%%%%%%%%%%%%%%%%%%%%%%%%

The exponential distribution satisfies

\[
\boxed{
P
(
T>s+t
\mid
T>s
)
=
P(T>t).
}
\]

Thus surviving to time \(s\) does not change the distribution of the
remaining lifetime.

This memoryless property is equivalent to a constant hazard.

%%%%%%%%%%%%%%%%%%%%%%%%%%%%%%%%%%%%%%%%%%%%%%%%%%%%%%%%%%%%
\subsection{Mean and Median Exponential Survival}
\label{subsec:exponential-mean-median}
%%%%%%%%%%%%%%%%%%%%%%%%%%%%%%%%%%%%%%%%%%%%%%%%%%%%%%%%%%%%

For

\[
T
\sim
\operatorname{Exponential}(\lambda),
\]

the mean is

\[
\boxed{
\mathbb{E}[T]
=
\frac1\lambda.
}
\]

The median satisfies

\[
e^{-\lambda m}
=
\frac12,
\]

so

\[
\boxed{
m
=
\frac{
\log2
}{
\lambda
}.
}
\]

%%%%%%%%%%%%%%%%%%%%%%%%%%%%%%%%%%%%%%%%%%%%%%%%%%%%%%%%%%%%
\subsection{Worked Example: Exponential Median Survival}
\label{subsec:worked-exponential-median}
%%%%%%%%%%%%%%%%%%%%%%%%%%%%%%%%%%%%%%%%%%%%%%%%%%%%%%%%%%%%

\begin{workedexamplebox}{Constant-Hazard Median}

Suppose

\[
\lambda=0.1.
\]

Then

\[
m
=
\frac{
\log2
}{
0.1
}.
\]

Therefore,

\begin{answerbox}
\[
\boxed{
m
\approx
6.93
}
\]

time units.
\end{answerbox}

\end{workedexamplebox}

%%%%%%%%%%%%%%%%%%%%%%%%%%%%%%%%%%%%%%%%%%%%%%%%%%%%%%%%%%%%
\subsection{Exponential Likelihood with Censoring}
\label{subsec:exponential-censored-likelihood}
%%%%%%%%%%%%%%%%%%%%%%%%%%%%%%%%%%%%%%%%%%%%%%%%%%%%%%%%%%%%

For right-censored observations,

\[
f(y_i)
=
\lambda e^{-\lambda y_i}
\]

when \(\delta_i=1\), and

\[
S(y_i)
=
e^{-\lambda y_i}
\]

when \(\delta_i=0\).

Therefore,

\[
L(\lambda)
=
\prod_i
\lambda^{\delta_i}
e^{-\lambda y_i}.
\]

Hence,

\[
\boxed{
L(\lambda)
=
\lambda^D
e^{-\lambda Y},
}
\]

where

\[
D
=
\sum_i\delta_i
\]

and

\[
Y
=
\sum_iy_i.
\]

%%%%%%%%%%%%%%%%%%%%%%%%%%%%%%%%%%%%%%%%%%%%%%%%%%%%%%%%%%%%
\subsection{Exponential MLE with Censoring}
\label{subsec:exponential-mle-censoring}
%%%%%%%%%%%%%%%%%%%%%%%%%%%%%%%%%%%%%%%%%%%%%%%%%%%%%%%%%%%%

The log-likelihood is

\[
\ell(\lambda)
=
D\log\lambda
-
\lambda Y.
\]

Differentiating,

\[
\ell'(\lambda)
=
\frac{D}{\lambda}
-
Y.
\]

Thus,

\[
\boxed{
\widehat{\lambda}
=
\frac{
D
}{
Y
}.
}
\]

This is:

\[
\boxed{
\frac{
\text{number of observed events}
}{
\text{total observed time at risk}
}.
}
\]

%%%%%%%%%%%%%%%%%%%%%%%%%%%%%%%%%%%%%%%%%%%%%%%%%%%%%%%%%%%%
\subsection{Weibull Distribution}
\label{subsec:weibull-survival}
%%%%%%%%%%%%%%%%%%%%%%%%%%%%%%%%%%%%%%%%%%%%%%%%%%%%%%%%%%%%

A Weibull model can be parameterized using shape \(k>0\) and scale
\(\lambda>0\) with survival function

\[
\boxed{
S(t)
=
\exp
\left[
-
\left(
\frac{t}{\lambda}
\right)^k
\right].
}
\]

The cumulative hazard is

\[
\boxed{
H(t)
=
\left(
\frac{t}{\lambda}
\right)^k.
}
\]

%%%%%%%%%%%%%%%%%%%%%%%%%%%%%%%%%%%%%%%%%%%%%%%%%%%%%%%%%%%%
\subsection{Weibull Hazard}
\label{subsec:weibull-hazard}
%%%%%%%%%%%%%%%%%%%%%%%%%%%%%%%%%%%%%%%%%%%%%%%%%%%%%%%%%%%%

The Weibull hazard is

\[
\boxed{
h(t)
=
\frac{k}{\lambda}
\left(
\frac{t}{\lambda}
\right)^{k-1}.
}
\]

Therefore:

\[
\boxed{
k=1
\rightarrow
\text{constant hazard},
}
\]

\[
\boxed{
k>1
\rightarrow
\text{increasing hazard},
}
\]

and

\[
\boxed{
0<k<1
\rightarrow
\text{decreasing hazard}.
}
\]

The exponential model is the special case

\[
\boxed{
k=1.
}
\]

%%%%%%%%%%%%%%%%%%%%%%%%%%%%%%%%%%%%%%%%%%%%%%%%%%%%%%%%%%%%
\subsection{Log-Normal Survival Model}
\label{subsec:lognormal-survival}
%%%%%%%%%%%%%%%%%%%%%%%%%%%%%%%%%%%%%%%%%%%%%%%%%%%%%%%%%%%%

A log-normal survival model assumes

\[
\boxed{
\log T
\sim
N(\mu,\sigma^2).
}
\]

Thus \(T\) is positive but its hazard need not be monotone.

The survival function is

\[
\boxed{
S(t)
=
1
-
\Phi
\left(
\frac{
\log t-\mu
}{
\sigma
}
\right).
}
\]

%%%%%%%%%%%%%%%%%%%%%%%%%%%%%%%%%%%%%%%%%%%%%%%%%%%%%%%%%%%%
\subsection{Log-Logistic Survival Model}
\label{subsec:loglogistic-survival}
%%%%%%%%%%%%%%%%%%%%%%%%%%%%%%%%%%%%%%%%%%%%%%%%%%%%%%%%%%%%

A log-logistic model arises when

\[
\log T
\]

has a logistic distribution.

Depending on parameters, its hazard may rise and then fall.

This can model event processes with nonmonotone hazard behavior.

%%%%%%%%%%%%%%%%%%%%%%%%%%%%%%%%%%%%%%%%%%%%%%%%%%%%%%%%%%%%
\subsection{Regression for Survival Data}
\label{subsec:survival-regression}
%%%%%%%%%%%%%%%%%%%%%%%%%%%%%%%%%%%%%%%%%%%%%%%%%%%%%%%%%%%%

Survival regression relates event-time distributions to predictors such
as:

\[
\boxed{
\text{treatment},
}
\]

\[
\boxed{
\text{age},
}
\]

\[
\boxed{
\text{exposure},
}
\]

or

\[
\boxed{
\text{machine conditions}.
}
\]

Two major frameworks are:

\[
\boxed{
\text{proportional hazards},
}
\]

and

\[
\boxed{
\text{accelerated failure time}.
}
\]

%%%%%%%%%%%%%%%%%%%%%%%%%%%%%%%%%%%%%%%%%%%%%%%%%%%%%%%%%%%%
\subsection{Proportional-Hazards Model}
\label{subsec:proportional-hazards}
%%%%%%%%%%%%%%%%%%%%%%%%%%%%%%%%%%%%%%%%%%%%%%%%%%%%%%%%%%%%

A proportional-hazards model assumes

\[
\boxed{
h
(
t\mid\mathbf{x}
)
=
h_0(t)
\exp
(
\mathbf{x}^T
\boldsymbol{\beta}
).
}
\]

Here,

\[
h_0(t)
\]

is the baseline hazard.

The covariates multiply the baseline hazard by a constant factor.

%%%%%%%%%%%%%%%%%%%%%%%%%%%%%%%%%%%%%%%%%%%%%%%%%%%%%%%%%%%%
\subsection{Cox Proportional-Hazards Model}
\label{subsec:cox-proportional-hazards}
%%%%%%%%%%%%%%%%%%%%%%%%%%%%%%%%%%%%%%%%%%%%%%%%%%%%%%%%%%%%

The Cox model is

\[
\boxed{
h
(
t\mid\mathbf{x}
)
=
h_0(t)
\exp
(
\beta_1x_1
+
\cdots
+
\beta_px_p
).
}
\]

The baseline hazard

\[
h_0(t)
\]

is left unspecified.

This makes the model semiparametric.

%%%%%%%%%%%%%%%%%%%%%%%%%%%%%%%%%%%%%%%%%%%%%%%%%%%%%%%%%%%%
\subsection{Hazard Ratio}
\label{subsec:hazard-ratio}
%%%%%%%%%%%%%%%%%%%%%%%%%%%%%%%%%%%%%%%%%%%%%%%%%%%%%%%%%%%%

Consider two covariate vectors

\[
\mathbf{x}_1
\]

and

\[
\mathbf{x}_0.
\]

Their hazard ratio is

\[
\boxed{
HR
=
\frac{
h(t\mid\mathbf{x}_1)
}{
h(t\mid\mathbf{x}_0)
}
=
\exp
\left[
(
\mathbf{x}_1-\mathbf{x}_0
)^T
\boldsymbol{\beta}
\right].
}
\]

The baseline hazard cancels.

%%%%%%%%%%%%%%%%%%%%%%%%%%%%%%%%%%%%%%%%%%%%%%%%%%%%%%%%%%%%
\subsection{Interpretation of a Cox Coefficient}
\label{subsec:cox-coefficient-interpretation}
%%%%%%%%%%%%%%%%%%%%%%%%%%%%%%%%%%%%%%%%%%%%%%%%%%%%%%%%%%%%

For a one-unit increase in predictor \(x_j\), holding other modeled
variables fixed,

\[
\boxed{
HR
=
e^{\beta_j}.
}
\]

If

\[
e^{\beta_j}=2,
\]

the modeled instantaneous event rate is twice as large at every time
under the proportional-hazards assumption.

This does not mean survival time is halved.

%%%%%%%%%%%%%%%%%%%%%%%%%%%%%%%%%%%%%%%%%%%%%%%%%%%%%%%%%%%%
\subsection{Worked Example: Cox Hazard Ratio}
\label{subsec:worked-cox-hazard-ratio}
%%%%%%%%%%%%%%%%%%%%%%%%%%%%%%%%%%%%%%%%%%%%%%%%%%%%%%%%%%%%

\begin{workedexamplebox}{Interpreting a Cox Coefficient}

Suppose

\[
\widehat{\beta}=0.4.
\]

Then

\[
\widehat{HR}
=
e^{0.4}.
\]

Therefore,

\begin{answerbox}
\[
\boxed{
\widehat{HR}
\approx
1.49.
}
\]

A one-unit increase in the predictor is associated with approximately a
\(49\%\) higher modeled hazard, holding other modeled variables fixed.
\end{answerbox}

\end{workedexamplebox}

%%%%%%%%%%%%%%%%%%%%%%%%%%%%%%%%%%%%%%%%%%%%%%%%%%%%%%%%%%%%
\subsection{Why Cox Is Semiparametric}
\label{subsec:why-cox-semiparametric}
%%%%%%%%%%%%%%%%%%%%%%%%%%%%%%%%%%%%%%%%%%%%%%%%%%%%%%%%%%%%

The model specifies the finite-dimensional regression component

\[
\boldsymbol{\beta}
\]

but does not specify a finite-dimensional parametric family for

\[
h_0(t).
\]

Thus it contains both:

\[
\boxed{
\text{parametric regression coefficients}
}
\]

and

\[
\boxed{
\text{nonparametric baseline hazard}.
}
\]

%%%%%%%%%%%%%%%%%%%%%%%%%%%%%%%%%%%%%%%%%%%%%%%%%%%%%%%%%%%%
\subsection{Cox Partial Likelihood}
\label{subsec:cox-partial-likelihood}
%%%%%%%%%%%%%%%%%%%%%%%%%%%%%%%%%%%%%%%%%%%%%%%%%%%%%%%%%%%%

Suppose one event occurs at time \(t_j\) for subject \(i_j\).

Conditional on one event occurring among subjects in risk set \(R_j\),
the Cox model implies

\[
\boxed{
P
(
i_j
\text{ fails at }t_j
\mid
\text{one failure in }R_j
)
=
\frac{
\exp
(
\mathbf{x}_{i_j}^T
\boldsymbol{\beta}
)
}{
\sum_{k\in R_j}
\exp
(
\mathbf{x}_k^T
\boldsymbol{\beta}
)
}.
}
\]

Multiplying these terms gives the partial likelihood.

%%%%%%%%%%%%%%%%%%%%%%%%%%%%%%%%%%%%%%%%%%%%%%%%%%%%%%%%%%%%
\subsection{Partial Likelihood Formula}
\label{subsec:cox-partial-likelihood-formula}
%%%%%%%%%%%%%%%%%%%%%%%%%%%%%%%%%%%%%%%%%%%%%%%%%%%%%%%%%%%%

Ignoring tied event complications, the partial likelihood is

\[
\boxed{
L_P
(
\boldsymbol{\beta}
)
=
\prod_{j}
\frac{
\exp
(
\mathbf{x}_{i_j}^T
\boldsymbol{\beta}
)
}{
\displaystyle
\sum_{k\in R_j}
\exp
(
\mathbf{x}_k^T
\boldsymbol{\beta}
)
}.
}
\]

The baseline hazard cancels from each conditional probability.

Therefore \(\boldsymbol{\beta}\) can be estimated without specifying
\(h_0(t)\).

%%%%%%%%%%%%%%%%%%%%%%%%%%%%%%%%%%%%%%%%%%%%%%%%%%%%%%%%%%%%
\subsection{Tied Event Times}
\label{subsec:cox-ties}
%%%%%%%%%%%%%%%%%%%%%%%%%%%%%%%%%%%%%%%%%%%%%%%%%%%%%%%%%%%%

When several subjects experience events at the same recorded time, the
simple partial likelihood requires modification.

Common approximations include:

\[
\boxed{
\text{Breslow-type handling},
}
\]

and

\[
\boxed{
\text{Efron-type handling}.
}
\]

Exact approaches may also be used in certain settings.

%%%%%%%%%%%%%%%%%%%%%%%%%%%%%%%%%%%%%%%%%%%%%%%%%%%%%%%%%%%%
\subsection{Baseline Survival Under Cox Regression}
\label{subsec:cox-baseline-survival}
%%%%%%%%%%%%%%%%%%%%%%%%%%%%%%%%%%%%%%%%%%%%%%%%%%%%%%%%%%%%

The baseline cumulative hazard is

\[
\boxed{
H_0(t).
}
\]

Since

\[
h(t\mid\mathbf{x})
=
h_0(t)
e^{\mathbf{x}^T\boldsymbol{\beta}},
\]

the cumulative hazard is

\[
\boxed{
H
(
t\mid\mathbf{x}
)
=
H_0(t)
e^{\mathbf{x}^T\boldsymbol{\beta}}.
}
\]

Therefore,

\[
\boxed{
S
(
t\mid\mathbf{x}
)
=
[
S_0(t)
]^{
e^{\mathbf{x}^T\boldsymbol{\beta}}
}.
}
\]

%%%%%%%%%%%%%%%%%%%%%%%%%%%%%%%%%%%%%%%%%%%%%%%%%%%%%%%%%%%%
\subsection{Proportional-Hazards Assumption}
\label{subsec:ph-assumption}
%%%%%%%%%%%%%%%%%%%%%%%%%%%%%%%%%%%%%%%%%%%%%%%%%%%%%%%%%%%%

For two covariate patterns,

\[
\mathbf{x}_1
\]

and

\[
\mathbf{x}_0,
\]

the Cox model assumes

\[
\boxed{
\frac{
h(t\mid\mathbf{x}_1)
}{
h(t\mid\mathbf{x}_0)
}
}
\]

does not depend on \(t\).

Thus hazard ratios are constant over time.

This is the proportional-hazards assumption.

%%%%%%%%%%%%%%%%%%%%%%%%%%%%%%%%%%%%%%%%%%%%%%%%%%%%%%%%%%%%
\subsection{Nonproportional Hazards}
\label{subsec:nonproportional-hazards}
%%%%%%%%%%%%%%%%%%%%%%%%%%%%%%%%%%%%%%%%%%%%%%%%%%%%%%%%%%%%

If treatment effects change over time, the hazard ratio may not be
constant.

For example:

\[
\boxed{
HR(t)
}
\]

may be less than \(1\) early and greater than \(1\) later.

A standard Cox coefficient would then compress a time-varying effect into
a single summary.

%%%%%%%%%%%%%%%%%%%%%%%%%%%%%%%%%%%%%%%%%%%%%%%%%%%%%%%%%%%%
\subsection{Time-by-Covariate Interaction}
\label{subsec:time-covariate-interaction}
%%%%%%%%%%%%%%%%%%%%%%%%%%%%%%%%%%%%%%%%%%%%%%%%%%%%%%%%%%%%

One way to model a changing effect is

\[
\boxed{
h(t\mid x)
=
h_0(t)
\exp
[
\beta_1x
+
\beta_2xg(t)
].
}
\]

The hazard ratio for a one-unit change in \(x\) becomes

\[
\boxed{
HR(t)
=
\exp
[
\beta_1+\beta_2g(t)
].
}
\]

Thus the covariate effect can vary with time.

%%%%%%%%%%%%%%%%%%%%%%%%%%%%%%%%%%%%%%%%%%%%%%%%%%%%%%%%%%%%
\subsection{Schoenfeld Residuals Preview}
\label{subsec:schoenfeld-residuals-preview}
%%%%%%%%%%%%%%%%%%%%%%%%%%%%%%%%%%%%%%%%%%%%%%%%%%%%%%%%%%%%

Schoenfeld residuals are commonly used to investigate proportional
hazards for Cox-model covariates.

Systematic residual patterns against time may suggest time-varying
effects.

No diagnostic alone proves that the model is correct.

%%%%%%%%%%%%%%%%%%%%%%%%%%%%%%%%%%%%%%%%%%%%%%%%%%%%%%%%%%%%
\subsection{Stratified Cox Model}
\label{subsec:stratified-cox}
%%%%%%%%%%%%%%%%%%%%%%%%%%%%%%%%%%%%%%%%%%%%%%%%%%%%%%%%%%%%

A stratified Cox model allows different baseline hazards for strata:

\[
\boxed{
h_k
(
t\mid\mathbf{x}
)
=
h_{0k}(t)
\exp
(
\mathbf{x}^T
\boldsymbol{\beta}
).
}
\]

The regression coefficients are shared, but the baseline hazard can
differ across strata.

This is useful for controlling a categorical variable whose
proportional-hazards effect is not of direct interest.

%%%%%%%%%%%%%%%%%%%%%%%%%%%%%%%%%%%%%%%%%%%%%%%%%%%%%%%%%%%%
\subsection{Time-Dependent Covariates}
\label{subsec:time-dependent-covariates}
%%%%%%%%%%%%%%%%%%%%%%%%%%%%%%%%%%%%%%%%%%%%%%%%%%%%%%%%%%%%

Some predictors change during follow-up.

Write

\[
\boxed{
\mathbf{x}(t).
}
\]

Then a Cox model may take the form

\[
\boxed{
h
(
t\mid\mathbf{x}(t)
)
=
h_0(t)
\exp
[
\mathbf{x}(t)^T
\boldsymbol{\beta}
].
}
\]

Risk sets must use predictor values relevant at each event time.

%%%%%%%%%%%%%%%%%%%%%%%%%%%%%%%%%%%%%%%%%%%%%%%%%%%%%%%%%%%%
\subsection{Internal and External Time-Dependent Covariates}
\label{subsec:internal-external-time-covariates}
%%%%%%%%%%%%%%%%%%%%%%%%%%%%%%%%%%%%%%%%%%%%%%%%%%%%%%%%%%%%

An external time-dependent covariate evolves independently of the
subject's event process in a suitable sense.

Examples may include calendar-level environmental exposures.

An internal covariate is generated by the subject's evolving health or
state process and can require more careful interpretation.

%%%%%%%%%%%%%%%%%%%%%%%%%%%%%%%%%%%%%%%%%%%%%%%%%%%%%%%%%%%%
\subsection{Accelerated Failure-Time Models}
\label{subsec:accelerated-failure-time}
%%%%%%%%%%%%%%%%%%%%%%%%%%%%%%%%%%%%%%%%%%%%%%%%%%%%%%%%%%%%

Accelerated failure-time, or AFT, models act directly on survival time.

A common form is

\[
\boxed{
\log T
=
\mathbf{x}^T
\boldsymbol{\beta}
+
\varepsilon.
}
\]

Exponentiating,

\[
\boxed{
T
=
e^{\mathbf{x}^T\boldsymbol{\beta}}
e^\varepsilon.
}
\]

Covariates multiply the time scale.

%%%%%%%%%%%%%%%%%%%%%%%%%%%%%%%%%%%%%%%%%%%%%%%%%%%%%%%%%%%%
\subsection{Time Ratio}
\label{subsec:time-ratio}
%%%%%%%%%%%%%%%%%%%%%%%%%%%%%%%%%%%%%%%%%%%%%%%%%%%%%%%%%%%%

For a one-unit increase in \(x_j\), an AFT model multiplies survival time
by

\[
\boxed{
e^{\beta_j}.
}
\]

This is a time ratio.

If

\[
e^{\beta_j}=1.5,
\]

the modeled event time is stretched by a factor of \(1.5\).

This interpretation differs from a hazard ratio.

%%%%%%%%%%%%%%%%%%%%%%%%%%%%%%%%%%%%%%%%%%%%%%%%%%%%%%%%%%%%
\subsection{Hazard Ratio Versus Time Ratio}
\label{subsec:hazard-ratio-versus-time-ratio}
%%%%%%%%%%%%%%%%%%%%%%%%%%%%%%%%%%%%%%%%%%%%%%%%%%%%%%%%%%%%

The Cox model uses

\[
\boxed{
\text{hazard ratios}.
}
\]

AFT models use

\[
\boxed{
\text{time ratios}.
}
\]

These describe different features of the survival distribution and should
not be interpreted interchangeably.

%%%%%%%%%%%%%%%%%%%%%%%%%%%%%%%%%%%%%%%%%%%%%%%%%%%%%%%%%%%%
\subsection{Parametric AFT Models}
\label{subsec:parametric-aft-models}
%%%%%%%%%%%%%%%%%%%%%%%%%%%%%%%%%%%%%%%%%%%%%%%%%%%%%%%%%%%%

AFT models arise naturally from distributions for

\[
\log T.
\]

Examples include:

\[
\boxed{
\text{log-normal AFT},
}
\]

\[
\boxed{
\text{log-logistic AFT},
}
\]

and certain Weibull parameterizations.

Parametric AFT models can support direct extrapolation beyond observed
follow-up, but extrapolation relies strongly on the assumed distribution.

%%%%%%%%%%%%%%%%%%%%%%%%%%%%%%%%%%%%%%%%%%%%%%%%%%%%%%%%%%%%
\subsection{Survival Regression Diagnostics}
\label{subsec:survival-regression-diagnostics}
%%%%%%%%%%%%%%%%%%%%%%%%%%%%%%%%%%%%%%%%%%%%%%%%%%%%%%%%%%%%

Useful diagnostics include:

\[
\boxed{
\text{proportional-hazards checks},
}
\]

\[
\boxed{
\text{functional-form checks},
}
\]

\[
\boxed{
\text{influential-observation checks},
}
\]

and

\[
\boxed{
\text{residual analysis}.
}
\]

Model assumptions should be assessed rather than inferred solely from
goodness-of-fit statistics.

%%%%%%%%%%%%%%%%%%%%%%%%%%%%%%%%%%%%%%%%%%%%%%%%%%%%%%%%%%%%
\subsection{Martingale Residuals Preview}
\label{subsec:martingale-residuals-preview}
%%%%%%%%%%%%%%%%%%%%%%%%%%%%%%%%%%%%%%%%%%%%%%%%%%%%%%%%%%%%

Martingale residuals compare observed event counts with estimated
cumulative event intensity.

They can help assess nonlinear covariate effects and unusual
observations.

Their distribution is generally skewed.

%%%%%%%%%%%%%%%%%%%%%%%%%%%%%%%%%%%%%%%%%%%%%%%%%%%%%%%%%%%%
\subsection{Deviance Residuals Preview}
\label{subsec:survival-deviance-residuals}
%%%%%%%%%%%%%%%%%%%%%%%%%%%%%%%%%%%%%%%%%%%%%%%%%%%%%%%%%%%%

Deviance residuals transform martingale residuals to produce a more
symmetric diagnostic scale.

They are useful for identifying observations that fit the Cox model
poorly.

%%%%%%%%%%%%%%%%%%%%%%%%%%%%%%%%%%%%%%%%%%%%%%%%%%%%%%%%%%%%
\subsection{Competing Risks}
\label{subsec:competing-risks}
%%%%%%%%%%%%%%%%%%%%%%%%%%%%%%%%%%%%%%%%%%%%%%%%%%%%%%%%%%%%

In some studies, several event types can occur and one event can prevent
observation of another.

Examples include different causes of failure or death.

Let

\[
J
\]

denote the event type.

Then the outcome includes both:

\[
\boxed{
T
}
\]

and

\[
\boxed{
J.
}
\]

%%%%%%%%%%%%%%%%%%%%%%%%%%%%%%%%%%%%%%%%%%%%%%%%%%%%%%%%%%%%
\subsection{Cause-Specific Hazard}
\label{subsec:cause-specific-hazard}
%%%%%%%%%%%%%%%%%%%%%%%%%%%%%%%%%%%%%%%%%%%%%%%%%%%%%%%%%%%%

For cause \(k\), the cause-specific hazard is

\[
\boxed{
h_k(t)
=
\lim_{\Delta t\to0^+}
\frac{
P
(
t\leq T<t+\Delta t,
J=k
\mid
T\geq t
)
}{
\Delta t
}.
}
\]

The overall hazard is

\[
\boxed{
h(t)
=
\sum_k
h_k(t).
}
\]

%%%%%%%%%%%%%%%%%%%%%%%%%%%%%%%%%%%%%%%%%%%%%%%%%%%%%%%%%%%%
\subsection{Cumulative Incidence Function}
\label{subsec:cumulative-incidence-function}
%%%%%%%%%%%%%%%%%%%%%%%%%%%%%%%%%%%%%%%%%%%%%%%%%%%%%%%%%%%%

The cumulative incidence function for cause \(k\) is

\[
\boxed{
F_k(t)
=
P
(
T\leq t,
J=k
).
}
\]

It gives the probability of experiencing cause \(k\) by time \(t\) in
the presence of competing events.

%%%%%%%%%%%%%%%%%%%%%%%%%%%%%%%%%%%%%%%%%%%%%%%%%%%%%%%%%%%%
\subsection{Cause-Specific Hazard and Cumulative Incidence}
\label{subsec:csh-versus-cif}
%%%%%%%%%%%%%%%%%%%%%%%%%%%%%%%%%%%%%%%%%%%%%%%%%%%%%%%%%%%%

Cause-specific hazards describe instantaneous cause-specific rates among
individuals still event-free.

Cumulative incidence describes the actual probability of observing a
specific event type over time.

They answer different questions.

A covariate can affect cause-specific hazards and cumulative incidence in
different ways because competing events also influence probability.

%%%%%%%%%%%%%%%%%%%%%%%%%%%%%%%%%%%%%%%%%%%%%%%%%%%%%%%%%%%%
\subsection{Why Kaplan--Meier Can Fail for Competing Risks}
\label{subsec:km-competing-risks}
%%%%%%%%%%%%%%%%%%%%%%%%%%%%%%%%%%%%%%%%%%%%%%%%%%%%%%%%%%%%

Treating competing events as ordinary censoring and using

\[
1-\widehat{S}(t)
\]

can overestimate cause-specific cumulative incidence.

A competing event changes whether the event of interest can still occur.

Therefore cumulative incidence methods should be used when event type is
important.

%%%%%%%%%%%%%%%%%%%%%%%%%%%%%%%%%%%%%%%%%%%%%%%%%%%%%%%%%%%%
\subsection{Subdistribution Hazard Preview}
\label{subsec:subdistribution-hazard-preview}
%%%%%%%%%%%%%%%%%%%%%%%%%%%%%%%%%%%%%%%%%%%%%%%%%%%%%%%%%%%%

Another competing-risks modeling approach uses a subdistribution hazard
constructed to model cumulative incidence directly.

This leads to regression methods distinct from cause-specific Cox
regression.

The interpretation should be stated carefully because the risk sets are
defined differently.

%%%%%%%%%%%%%%%%%%%%%%%%%%%%%%%%%%%%%%%%%%%%%%%%%%%%%%%%%%%%
\subsection{Recurrent Events}
\label{subsec:recurrent-events}
%%%%%%%%%%%%%%%%%%%%%%%%%%%%%%%%%%%%%%%%%%%%%%%%%%%%%%%%%%%%

Some subjects can experience the event multiple times.

Examples include:

\[
\boxed{
\text{hospital admissions},
}
\]

\[
\boxed{
\text{equipment breakdowns},
}
\]

or

\[
\boxed{
\text{disease relapses}.
}
\]

These recurrent event times within a subject are dependent.

%%%%%%%%%%%%%%%%%%%%%%%%%%%%%%%%%%%%%%%%%%%%%%%%%%%%%%%%%%%%
\subsection{Counting-Process Representation}
\label{subsec:counting-process-survival}
%%%%%%%%%%%%%%%%%%%%%%%%%%%%%%%%%%%%%%%%%%%%%%%%%%%%%%%%%%%%

Let

\[
N_i(t)
\]

be the number of observed events for subject \(i\) by time \(t\).

Then recurrent event data can be represented as counting processes.

The at-risk indicator may be written

\[
\boxed{
Y_i(t).
}
\]

This notation forms the basis of advanced modern survival theory.

%%%%%%%%%%%%%%%%%%%%%%%%%%%%%%%%%%%%%%%%%%%%%%%%%%%%%%%%%%%%
\subsection{Intensity Process Preview}
\label{subsec:survival-intensity-process}
%%%%%%%%%%%%%%%%%%%%%%%%%%%%%%%%%%%%%%%%%%%%%%%%%%%%%%%%%%%%

A counting process may have an intensity

\[
\boxed{
\lambda_i(t)
=
Y_i(t)
h_i(t).
}
\]

It represents the instantaneous expected event rate conditional on the
past.

This connects survival analysis directly with stochastic-process theory.

%%%%%%%%%%%%%%%%%%%%%%%%%%%%%%%%%%%%%%%%%%%%%%%%%%%%%%%%%%%%
\subsection{Frailty Models}
\label{subsec:frailty-models}
%%%%%%%%%%%%%%%%%%%%%%%%%%%%%%%%%%%%%%%%%%%%%%%%%%%%%%%%%%%%

A frailty model introduces an unobserved random effect representing
subject- or cluster-level risk.

A multiplicative frailty model may be

\[
\boxed{
h_i(t)
=
Z_i
h_0(t)
\exp
(
\mathbf{x}_i^T
\boldsymbol{\beta}
),
}
\]

where

\[
Z_i>0
\]

is an unobserved frailty.

%%%%%%%%%%%%%%%%%%%%%%%%%%%%%%%%%%%%%%%%%%%%%%%%%%%%%%%%%%%%
\subsection{Interpretation of Frailty}
\label{subsec:frailty-interpretation}
%%%%%%%%%%%%%%%%%%%%%%%%%%%%%%%%%%%%%%%%%%%%%%%%%%%%%%%%%%%%

Individuals with larger frailty values have larger hazards.

Frailty models can account for:

\[
\boxed{
\text{unobserved heterogeneity},
}
\]

or

\[
\boxed{
\text{dependence within clusters}.
}
\]

They are analogous to random-effects models in longitudinal and
hierarchical data analysis.

%%%%%%%%%%%%%%%%%%%%%%%%%%%%%%%%%%%%%%%%%%%%%%%%%%%%%%%%%%%%
\subsection{Shared Frailty}
\label{subsec:shared-frailty}
%%%%%%%%%%%%%%%%%%%%%%%%%%%%%%%%%%%%%%%%%%%%%%%%%%%%%%%%%%%%

Subjects within a common cluster may share the same frailty:

\[
\boxed{
Z_g.
}
\]

For example, members of the same family, hospital, or machine system may
share unmeasured risk factors.

This induces positive within-cluster event dependence.

%%%%%%%%%%%%%%%%%%%%%%%%%%%%%%%%%%%%%%%%%%%%%%%%%%%%%%%%%%%%
\subsection{Discrete-Time Survival}
\label{subsec:discrete-time-survival}
%%%%%%%%%%%%%%%%%%%%%%%%%%%%%%%%%%%%%%%%%%%%%%%%%%%%%%%%%%%%

If events can occur only at discrete time points

\[
t=1,2,\ldots,
\]

define the discrete hazard

\[
\boxed{
h_t
=
P
(
T=t
\mid
T\geq t
).
}
\]

Then

\[
\boxed{
S(t)
=
\prod_{j=1}^{t}
(
1-h_j
).
}
\]

Discrete-time survival can be modeled using binary regression on
person-period data.

%%%%%%%%%%%%%%%%%%%%%%%%%%%%%%%%%%%%%%%%%%%%%%%%%%%%%%%%%%%%
\subsection{Complementary Log-Log Link Preview}
\label{subsec:cloglog-survival}
%%%%%%%%%%%%%%%%%%%%%%%%%%%%%%%%%%%%%%%%%%%%%%%%%%%%%%%%%%%%

A common link for discrete hazards is the complementary log-log link:

\[
\boxed{
\log
[
-\log(1-h_t)
]
=
\alpha_t
+
\mathbf{x}^T
\boldsymbol{\beta}.
}
\]

This has a close connection with grouped proportional-hazards models.

%%%%%%%%%%%%%%%%%%%%%%%%%%%%%%%%%%%%%%%%%%%%%%%%%%%%%%%%%%%%
\subsection{Reliability Function}
\label{subsec:reliability-function}
%%%%%%%%%%%%%%%%%%%%%%%%%%%%%%%%%%%%%%%%%%%%%%%%%%%%%%%%%%%%

In engineering, the survival function is often called the reliability
function:

\[
\boxed{
R(t)
=
P(T>t).
}
\]

Thus

\[
\boxed{
R(t)=S(t).
}
\]

The hazard may be called the failure rate.

%%%%%%%%%%%%%%%%%%%%%%%%%%%%%%%%%%%%%%%%%%%%%%%%%%%%%%%%%%%%
\subsection{Bathtub Hazard}
\label{subsec:bathtub-hazard}
%%%%%%%%%%%%%%%%%%%%%%%%%%%%%%%%%%%%%%%%%%%%%%%%%%%%%%%%%%%%

Engineering systems often display a bathtub-shaped hazard:

\[
\boxed{
\text{high early failure rate}
}
\]

followed by

\[
\boxed{
\text{low stable failure rate}
}
\]

and eventually

\[
\boxed{
\text{increasing wear-out failure rate}.
}
\]

Simple exponential models cannot represent this shape.

%%%%%%%%%%%%%%%%%%%%%%%%%%%%%%%%%%%%%%%%%%%%%%%%%%%%%%%%%%%%
\subsection{Restricted Mean Survival Time}
\label{subsec:restricted-mean-survival-time}
%%%%%%%%%%%%%%%%%%%%%%%%%%%%%%%%%%%%%%%%%%%%%%%%%%%%%%%%%%%%

For a fixed horizon \(\tau\), the restricted mean survival time is

\[
\boxed{
RMST(\tau)
=
\mathbb{E}
[
\min(T,\tau)
].
}
\]

It can be expressed as

\[
\boxed{
RMST(\tau)
=
\int_0^\tau
S(t)\,dt.
}
\]

This summarizes expected event-free time up to a prespecified horizon.

%%%%%%%%%%%%%%%%%%%%%%%%%%%%%%%%%%%%%%%%%%%%%%%%%%%%%%%%%%%%
\subsection{Difference in Restricted Mean Survival}
\label{subsec:rmst-difference}
%%%%%%%%%%%%%%%%%%%%%%%%%%%%%%%%%%%%%%%%%%%%%%%%%%%%%%%%%%%%

For two groups,

\[
\boxed{
\Delta_{RMST}
=
\int_0^\tau
[
S_1(t)-S_0(t)
]
\,dt.
}
\]

This gives a time-scale treatment effect.

It can be useful when proportional hazards is questionable.

%%%%%%%%%%%%%%%%%%%%%%%%%%%%%%%%%%%%%%%%%%%%%%%%%%%%%%%%%%%%
\subsection{Hazard Ratio Limitations}
\label{subsec:hazard-ratio-limitations}
%%%%%%%%%%%%%%%%%%%%%%%%%%%%%%%%%%%%%%%%%%%%%%%%%%%%%%%%%%%%

A hazard ratio is not:

\[
\boxed{
\text{a risk ratio},
}
\]

\[
\boxed{
\text{a survival-probability ratio},
}
\]

or

\[
\boxed{
\text{a ratio of median survival times}.
}
\]

Its interpretation is conditional on remaining event-free and relies on
the model being used.

%%%%%%%%%%%%%%%%%%%%%%%%%%%%%%%%%%%%%%%%%%%%%%%%%%%%%%%%%%%%
\subsection{Survival Probability at a Fixed Time}
\label{subsec:survival-fixed-time}
%%%%%%%%%%%%%%%%%%%%%%%%%%%%%%%%%%%%%%%%%%%%%%%%%%%%%%%%%%%%

Sometimes the most interpretable quantity is

\[
\boxed{
S(t^*)
}
\]

at a clinically or scientifically meaningful time \(t^*\).

Two groups may be compared through

\[
\boxed{
S_1(t^*)-S_0(t^*)
}
\]

or

\[
\boxed{
\frac{
S_1(t^*)
}{
S_0(t^*)
}.
}
\]

These differ conceptually from hazard ratios.

%%%%%%%%%%%%%%%%%%%%%%%%%%%%%%%%%%%%%%%%%%%%%%%%%%%%%%%%%%%%
\subsection{Landmark Analysis Preview}
\label{subsec:landmark-analysis}
%%%%%%%%%%%%%%%%%%%%%%%%%%%%%%%%%%%%%%%%%%%%%%%%%%%%%%%%%%%%

A landmark analysis selects a prespecified time \(t_L\) and analyzes
subsequent outcomes among subjects still at risk at \(t_L\).

This can help address time-dependent predictors when the landmark
definition is scientifically meaningful.

Selection at the landmark changes the target population and must be
interpreted accordingly.

%%%%%%%%%%%%%%%%%%%%%%%%%%%%%%%%%%%%%%%%%%%%%%%%%%%%%%%%%%%%
\subsection{Immortal-Time Bias Preview}
\label{subsec:immortal-time-bias}
%%%%%%%%%%%%%%%%%%%%%%%%%%%%%%%%%%%%%%%%%%%%%%%%%%%%%%%%%%%%

Immortal-time bias occurs when exposure classification requires subjects
to survive for some period before they can enter an exposed group.

If that guaranteed survival time is incorrectly attributed to the
exposure group, treatment or exposure effects can be severely biased.

Time-dependent exposure modeling is often required.

%%%%%%%%%%%%%%%%%%%%%%%%%%%%%%%%%%%%%%%%%%%%%%%%%%%%%%%%%%%%
\subsection{Left Truncation Risk Sets}
\label{subsec:left-truncation-risk-sets}
%%%%%%%%%%%%%%%%%%%%%%%%%%%%%%%%%%%%%%%%%%%%%%%%%%%%%%%%%%%%

With delayed entry, subject \(i\) enters the risk set only when

\[
\boxed{
L_i<t\leq Y_i.
}
\]

Subjects who have not yet entered must not contribute to the risk set.

This is essential for valid Kaplan--Meier and Cox analyses under left
truncation.

%%%%%%%%%%%%%%%%%%%%%%%%%%%%%%%%%%%%%%%%%%%%%%%%%%%%%%%%%%%%
\subsection{Interval-Censored Survival Estimation Preview}
\label{subsec:interval-censored-estimation}
%%%%%%%%%%%%%%%%%%%%%%%%%%%%%%%%%%%%%%%%%%%%%%%%%%%%%%%%%%%%

Ordinary Kaplan--Meier methods assume exact event times or right
censoring.

For interval-censored data, the event time is known only to lie in

\[
(L_i,R_i].
\]

Special nonparametric or parametric likelihood methods are required.

%%%%%%%%%%%%%%%%%%%%%%%%%%%%%%%%%%%%%%%%%%%%%%%%%%%%%%%%%%%%
\subsection{Survival Prediction}
\label{subsec:survival-prediction}
%%%%%%%%%%%%%%%%%%%%%%%%%%%%%%%%%%%%%%%%%%%%%%%%%%%%%%%%%%%%

A survival model may predict an entire function

\[
\boxed{
\widehat{S}(t\mid\mathbf{x}).
}
\]

This differs from ordinary regression, where prediction often targets a
single scalar mean.

Useful predictions include:

\[
\boxed{
\text{survival probability at }t,
}
\]

\[
\boxed{
\text{median survival},
}
\]

and

\[
\boxed{
\text{cumulative event risk}.
}
\]

%%%%%%%%%%%%%%%%%%%%%%%%%%%%%%%%%%%%%%%%%%%%%%%%%%%%%%%%%%%%
\subsection{Censoring and Predictive Evaluation}
\label{subsec:censoring-predictive-evaluation}
%%%%%%%%%%%%%%%%%%%%%%%%%%%%%%%%%%%%%%%%%%%%%%%%%%%%%%%%%%%%

Prediction accuracy cannot be evaluated naively when outcomes are
censored.

For example, at a fixed prediction horizon some subjects may have been
censored before their event status is known.

Evaluation methods may therefore use:

\[
\boxed{
\text{risk-set methods},
}
\]

\[
\boxed{
\text{inverse probability of censoring weights},
}
\]

or specialized survival scoring rules.

%%%%%%%%%%%%%%%%%%%%%%%%%%%%%%%%%%%%%%%%%%%%%%%%%%%%%%%%%%%%
\subsection{Concordance Index Preview}
\label{subsec:concordance-index-preview}
%%%%%%%%%%%%%%%%%%%%%%%%%%%%%%%%%%%%%%%%%%%%%%%%%%%%%%%%%%%%

The concordance index, or C-index, measures how well predicted risk
scores rank event times.

Conceptually, it estimates the probability that, for an informative pair,
the subject predicted to be at higher risk experiences the event earlier.

Censoring complicates which pairs are comparable.

%%%%%%%%%%%%%%%%%%%%%%%%%%%%%%%%%%%%%%%%%%%%%%%%%%%%%%%%%%%%
\subsection{Brier Score Preview}
\label{subsec:survival-brier-score}
%%%%%%%%%%%%%%%%%%%%%%%%%%%%%%%%%%%%%%%%%%%%%%%%%%%%%%%%%%%%

At prediction time \(t\), a survival Brier score compares predicted
survival probability with observed event status.

A simplified squared error is

\[
\boxed{
[
\mathbf{1}_{\{T>t\}}
-
\widehat{S}(t)
]^2.
}
\]

With censoring, weighting adjustments are required.

%%%%%%%%%%%%%%%%%%%%%%%%%%%%%%%%%%%%%%%%%%%%%%%%%%%%%%%%%%%%
\subsection{Calibration in Survival Models}
\label{subsec:survival-calibration}
%%%%%%%%%%%%%%%%%%%%%%%%%%%%%%%%%%%%%%%%%%%%%%%%%%%%%%%%%%%%

A survival prediction model is calibrated if predicted event
probabilities agree with observed event frequencies over time.

For example, among individuals with predicted

\[
80\%
\]

survival at time \(t\), approximately \(80\%\) should remain event-free
to that horizon in a well-calibrated population.

%%%%%%%%%%%%%%%%%%%%%%%%%%%%%%%%%%%%%%%%%%%%%%%%%%%%%%%%%%%%
\subsection{Survival Analysis Workflow}
\label{subsec:survival-analysis-workflow}
%%%%%%%%%%%%%%%%%%%%%%%%%%%%%%%%%%%%%%%%%%%%%%%%%%%%%%%%%%%%

A practical workflow is:

\begin{enumerate}
    \item define the event precisely;
    \item define the time origin;
    \item define the analysis time scale;
    \item identify censoring and truncation mechanisms;
    \item inspect follow-up distributions;
    \item summarize numbers at risk;
    \item estimate Kaplan--Meier curves;
    \item report survival probabilities at meaningful times;
    \item estimate median or restricted mean survival when appropriate;
    \item compare groups descriptively;
    \item use log-rank tests when their target is appropriate;
    \item fit regression models when covariate adjustment is needed;
    \item assess proportional hazards for Cox models;
    \item investigate nonlinear covariate effects;
    \item consider competing risks when multiple event types exist;
    \item consider recurrent-event or frailty methods for repeated or
          clustered outcomes;
    \item evaluate predictive discrimination and calibration when
          prediction is the goal;
    \item report assumptions and censoring information clearly.
\end{enumerate}

%%%%%%%%%%%%%%%%%%%%%%%%%%%%%%%%%%%%%%%%%%%%%%%%%%%%%%%%%%%%
\subsection{Choosing a Survival Method}
\label{subsec:choosing-survival-method}
%%%%%%%%%%%%%%%%%%%%%%%%%%%%%%%%%%%%%%%%%%%%%%%%%%%%%%%%%%%%

To estimate survival without covariates:

\[
\boxed{
\text{Kaplan--Meier}.
}
\]

To estimate cumulative hazard:

\[
\boxed{
\text{Nelson--Aalen}.
}
\]

To compare several survival curves:

\[
\boxed{
\text{log-rank test}.
}
\]

To model covariate hazard ratios:

\[
\boxed{
\text{Cox proportional hazards}.
}
\]

To model effects directly on survival time:

\[
\boxed{
\text{accelerated failure-time model}.
}
\]

For competing event types:

\[
\boxed{
\text{competing-risks methods}.
}
\]

%%%%%%%%%%%%%%%%%%%%%%%%%%%%%%%%%%%%%%%%%%%%%%%%%%%%%%%%%%%%
\subsection{Common Errors}
\label{subsec:survival-common-errors}
%%%%%%%%%%%%%%%%%%%%%%%%%%%%%%%%%%%%%%%%%%%%%%%%%%%%%%%%%%%%

\subsubsection{Treating Censored Times as Event Times}

A censoring time only establishes that the event has not occurred before
that time.

It is not the actual event time.

%%%%%%%%%%%%%%%%%%%%%%%%%%%%%%%%%%%%%%%%%%%%%%%%%%%%%%%%%%%%
\subsubsection{Deleting Censored Observations}

Censored observations still contribute information about survival.

Discarding them wastes information and can introduce bias.

%%%%%%%%%%%%%%%%%%%%%%%%%%%%%%%%%%%%%%%%%%%%%%%%%%%%%%%%%%%%
\subsubsection{Ignoring Informative Censoring}

Standard methods generally require censoring assumptions.

If subjects leave because of changing event risk, ordinary estimators may
be biased.

%%%%%%%%%%%%%%%%%%%%%%%%%%%%%%%%%%%%%%%%%%%%%%%%%%%%%%%%%%%%
\subsubsection{Confusing Censoring and Truncation}

Censored individuals appear in the dataset.

Truncated individuals may never enter the observed sample.

These require different risk-set and likelihood constructions.

%%%%%%%%%%%%%%%%%%%%%%%%%%%%%%%%%%%%%%%%%%%%%%%%%%%%%%%%%%%%
\subsubsection{Treating the Hazard as a Probability}

The hazard is an instantaneous rate conditional on survival.

It is not bounded above by \(1\).

%%%%%%%%%%%%%%%%%%%%%%%%%%%%%%%%%%%%%%%%%%%%%%%%%%%%%%%%%%%%
\subsubsection{Interpreting a Hazard Ratio as a Risk Ratio}

A hazard ratio compares instantaneous event rates among those still at
risk.

It does not equal a cumulative risk ratio in general.

%%%%%%%%%%%%%%%%%%%%%%%%%%%%%%%%%%%%%%%%%%%%%%%%%%%%%%%%%%%%
\subsubsection{Interpreting Hazard Ratio as a Time Ratio}

A Cox hazard ratio and an AFT time ratio refer to different model scales.

They should not be interchanged.

%%%%%%%%%%%%%%%%%%%%%%%%%%%%%%%%%%%%%%%%%%%%%%%%%%%%%%%%%%%%
\subsubsection{Assuming Proportional Hazards Without Checking}

The Cox model assumes covariate hazard ratios are constant over time.

Crossing or changing effects can violate this assumption.

%%%%%%%%%%%%%%%%%%%%%%%%%%%%%%%%%%%%%%%%%%%%%%%%%%%%%%%%%%%%
\subsubsection{Using Log-Rank Automatically When Curves Cross}

The ordinary log-rank test is optimized for certain persistent hazard
differences.

Strongly crossing hazards can produce low power or misleading summaries.

%%%%%%%%%%%%%%%%%%%%%%%%%%%%%%%%%%%%%%%%%%%%%%%%%%%%%%%%%%%%
\subsubsection{Reporting Only a \(p\)-Value}

Survival analyses should also report effect estimates and uncertainty,
such as:

\[
\boxed{
\text{hazard ratios},
}
\]

\[
\boxed{
\text{survival probabilities},
}
\]

or

\[
\boxed{
\text{RMST differences}.
}
\]

%%%%%%%%%%%%%%%%%%%%%%%%%%%%%%%%%%%%%%%%%%%%%%%%%%%%%%%%%%%%
\subsubsection{Ignoring Number at Risk}

Kaplan--Meier estimates near the end of follow-up may be based on very few
remaining subjects.

The number-at-risk table is essential for interpretation.

%%%%%%%%%%%%%%%%%%%%%%%%%%%%%%%%%%%%%%%%%%%%%%%%%%%%%%%%%%%%
\subsubsection{Extrapolating Kaplan--Meier Beyond Observed Follow-Up}

Kaplan--Meier is nonparametric and does not support meaningful
extrapolation beyond the available event information.

%%%%%%%%%%%%%%%%%%%%%%%%%%%%%%%%%%%%%%%%%%%%%%%%%%%%%%%%%%%%
\subsubsection{Assuming Parametric Extrapolation Is Automatically Reliable}

Parametric survival models can extrapolate, but the result depends
strongly on the assumed tail behavior.

Long-term predictions should be justified scientifically.

%%%%%%%%%%%%%%%%%%%%%%%%%%%%%%%%%%%%%%%%%%%%%%%%%%%%%%%%%%%%
\subsubsection{Using Kaplan--Meier for Cause-Specific Incidence with Competing Risks}

Treating competing events as ordinary censoring can overestimate
cause-specific cumulative incidence.

%%%%%%%%%%%%%%%%%%%%%%%%%%%%%%%%%%%%%%%%%%%%%%%%%%%%%%%%%%%%
\subsubsection{Ignoring Delayed Entry}

Subjects cannot contribute person-time before they become observable and
eligible for the risk set.

%%%%%%%%%%%%%%%%%%%%%%%%%%%%%%%%%%%%%%%%%%%%%%%%%%%%%%%%%%%%
\subsubsection{Creating Immortal-Time Bias}

Exposure status must be aligned correctly with time.

Future information cannot be used to classify earlier guaranteed
survival time as exposed.

%%%%%%%%%%%%%%%%%%%%%%%%%%%%%%%%%%%%%%%%%%%%%%%%%%%%%%%%%%%%
\subsubsection{Ignoring Within-Subject Dependence for Recurrent Events}

Multiple events from the same individual are correlated.

Ordinary independent-event analysis can underestimate uncertainty.

%%%%%%%%%%%%%%%%%%%%%%%%%%%%%%%%%%%%%%%%%%%%%%%%%%%%%%%%%%%%
\subsection{Applications}
\label{subsec:survival-applications}
%%%%%%%%%%%%%%%%%%%%%%%%%%%%%%%%%%%%%%%%%%%%%%%%%%%%%%%%%%%%

\begin{applicationbox}

\textbf{Medicine.}
Survival analysis models death, disease recurrence, treatment failure,
hospital discharge, and time to clinical events.

\medskip

\textbf{Engineering Reliability.}
Failure-time models estimate component lifetimes, hazard rates,
reliability functions, and maintenance schedules.

\medskip

\textbf{Business.}
Customer churn, employee retention, subscription duration, and time to
purchase can be modeled as event-time outcomes.

\medskip

\textbf{Manufacturing.}
Survival methods estimate time until machine breakdown, defect, or
replacement.

\medskip

\textbf{Epidemiology.}
Time from exposure to disease, recovery, hospitalization, or mortality
requires censoring-aware methods.

\medskip

\textbf{Education.}
Time to graduation, dropout, course completion, or program exit can be
treated as survival outcomes.

\medskip

\textbf{Finance.}
Time to default, bankruptcy, or contract termination can be analyzed
using hazard models.

\medskip

\textbf{Ecology.}
Survival of organisms, time to migration, and time to ecological events
can be studied with censored event-time models.

\end{applicationbox}

%%%%%%%%%%%%%%%%%%%%%%%%%%%%%%%%%%%%%%%%%%%%%%%%%%%%%%%%%%%%
\subsection{Exercises}
\label{subsec:survival-exercises}
%%%%%%%%%%%%%%%%%%%%%%%%%%%%%%%%%%%%%%%%%%%%%%%%%%%%%%%%%%%%

\begin{exercise}[1]
Define a survival time.
\end{exercise}

\begin{exercise}[1]
Define the survival function.
\end{exercise}

\begin{exercise}[1]
Define the hazard function.
\end{exercise}

\begin{exercise}[1]
Define the cumulative hazard function.
\end{exercise}

\begin{exercise}[1]
Define right censoring.
\end{exercise}

\begin{exercise}[1]
Define interval censoring.
\end{exercise}

\begin{exercise}[1]
Define left truncation.
\end{exercise}

\begin{exercise}[1]
Define a risk set.
\end{exercise}

\begin{exercise}[1]
State the Kaplan--Meier estimator.
\end{exercise}

\begin{exercise}[1]
State the Nelson--Aalen estimator.
\end{exercise}

\begin{exercise}[1]
Define a hazard ratio.
\end{exercise}

\begin{exercise}[1]
Define proportional hazards.
\end{exercise}

\begin{exercise}[1]
Define an accelerated failure-time model.
\end{exercise}

\begin{exercise}[1]
Define cumulative incidence.
\end{exercise}

\begin{exercise}[2]
Suppose

\[
S(t)=e^{-0.4t}.
\]

Find \(h(t)\).
\end{exercise}

\begin{exercise}[2]
For the same survival function, find \(H(t)\).
\end{exercise}

\begin{exercise}[2]
Suppose

\[
h(t)=3t^2.
\]

Find the cumulative hazard.
\end{exercise}

\begin{exercise}[2]
For the previous hazard, find the survival function.
\end{exercise}

\begin{exercise}[2]
For

\[
T\sim\operatorname{Exponential}(0.25),
\]

find the mean survival time.
\end{exercise}

\begin{exercise}[2]
For the previous distribution, find the median survival time.
\end{exercise}

\begin{exercise}[2]
Suppose the distinct event times have

\[
(n_1,d_1)=(8,1),
\]

and

\[
(n_2,d_2)=(6,2).
\]

Compute the Kaplan--Meier estimate after the second event time.
\end{exercise}

\begin{exercise}[2]
For the same data, compute the Nelson--Aalen cumulative hazard estimate.
\end{exercise}

\begin{exercise}[2]
Suppose

\[
\widehat{\beta}=0.7
\]

in a Cox model. Compute the corresponding hazard ratio.
\end{exercise}

\begin{exercise}[2]
Suppose an AFT coefficient is

\[
\widehat{\beta}=-0.2.
\]

Compute the corresponding time ratio.
\end{exercise}

\begin{exercise}[2]
Suppose a Weibull model has shape

\[
k=2.
\]

State whether its hazard increases, decreases, or remains constant.
\end{exercise}

\begin{exercise}[2]
Suppose a Weibull model has shape

\[
k=0.6.
\]

State the hazard behavior.
\end{exercise}

\begin{exercise}[3]
Derive

\[
h(t)
=
\frac{
f(t)
}{
S(t)
}.
\]
\end{exercise}

\begin{exercise}[3]
Derive

\[
H(t)
=
-\log S(t).
\]
\end{exercise}

\begin{exercise}[3]
Show that

\[
f(t)
=
h(t)e^{-H(t)}.
\]
\end{exercise}

\begin{exercise}[3]
Prove the tail-integral identity

\[
\mathbb{E}[T]
=
\int_0^\infty
S(t)\,dt
\]

for nonnegative continuous \(T\).
\end{exercise}

\begin{exercise}[3]
Derive the right-censored likelihood

\[
\prod_i
f(y_i)^{\delta_i}
S(y_i)^{1-\delta_i}.
\]
\end{exercise}

\begin{exercise}[3]
Derive the exponential MLE

\[
\widehat{\lambda}
=
\frac{
D
}{
Y
}
\]

under right censoring.
\end{exercise}

\begin{exercise}[3]
Explain why censoring changes risk sets but does not cause a Kaplan--Meier
drop.
\end{exercise}

\begin{exercise}[3]
Derive the Kaplan--Meier product-limit estimator from conditional
survival probabilities.
\end{exercise}

\begin{exercise}[3]
Explain the connection between Kaplan--Meier and Nelson--Aalen
estimators.
\end{exercise}

\begin{exercise}[3]
Derive the Weibull hazard from its survival function.
\end{exercise}

\begin{exercise}[3]
Show that the exponential distribution is a Weibull special case.
\end{exercise}

\begin{exercise}[3]
Prove the memoryless property of the exponential distribution.
\end{exercise}

\begin{exercise}[3]
Explain how the log-rank expected event count is constructed.
\end{exercise}

\begin{exercise}[3]
Explain why the Cox baseline hazard cancels from the partial likelihood.
\end{exercise}

\begin{exercise}[3]
Derive the Cox hazard ratio for two covariate vectors.
\end{exercise}

\begin{exercise}[3]
Derive

\[
S
(
t\mid\mathbf{x}
)
=
[
S_0(t)
]^{
e^{\mathbf{x}^T\boldsymbol{\beta}}
}.
\]
\end{exercise}

\begin{exercise}[3]
Explain the proportional-hazards assumption in words and mathematically.
\end{exercise}

\begin{exercise}[3]
Explain how a time-by-covariate interaction allows nonproportional
hazards.
\end{exercise}

\begin{exercise}[3]
Compare Cox and accelerated failure-time models.
\end{exercise}

\begin{exercise}[3]
Explain why a hazard ratio cannot generally be interpreted as a ratio of
median survival times.
\end{exercise}

\begin{exercise}[3]
Explain the difference between a cause-specific hazard and cumulative
incidence.
\end{exercise}

\begin{exercise}[3]
Explain why competing events should not always be treated as ordinary
censoring.
\end{exercise}

\begin{exercise}[3]
Explain why recurrent events within one individual are dependent.
\end{exercise}

\begin{exercise}[3]
Explain the purpose of a shared frailty model.
\end{exercise}

\begin{exercise}[3]
Derive the discrete survival relation

\[
S(t)
=
\prod_{j=1}^{t}
(
1-h_j
).
\]
\end{exercise}

\begin{exercise}[3]
Explain the meaning of restricted mean survival time.
\end{exercise}

\begin{exercise}[3]
Explain how immortal-time bias can arise.
\end{exercise}

\begin{exercise}[4]
Derive Greenwood's variance formula.
\end{exercise}

\begin{exercise}[4]
Study asymptotic properties of the Kaplan--Meier estimator.
\end{exercise}

\begin{exercise}[4]
Study martingale representations of Kaplan--Meier and Nelson--Aalen
estimators.
\end{exercise}

\begin{exercise}[4]
Derive the asymptotic null distribution of the log-rank statistic.
\end{exercise}

\begin{exercise}[4]
Study weighted log-rank tests.
\end{exercise}

\begin{exercise}[4]
Derive the Cox partial likelihood using conditional event
probabilities.
\end{exercise}

\begin{exercise}[4]
Derive the Cox score function and information matrix.
\end{exercise}

\begin{exercise}[4]
Study Breslow and Efron treatments of tied event times.
\end{exercise}

\begin{exercise}[4]
Investigate estimation of the baseline cumulative hazard in a Cox model.
\end{exercise}

\begin{exercise}[4]
Study Schoenfeld residuals and proportional-hazards diagnostics.
\end{exercise}

\begin{exercise}[4]
Study martingale and deviance residuals in Cox regression.
\end{exercise}

\begin{exercise}[4]
Derive likelihoods for parametric Weibull regression.
\end{exercise}

\begin{exercise}[4]
Study accelerated failure-time models under log-normal and log-logistic
errors.
\end{exercise}

\begin{exercise}[4]
Investigate competing-risks cumulative incidence estimation.
\end{exercise}

\begin{exercise}[4]
Study cause-specific Cox regression and subdistribution-hazard
regression.
\end{exercise}

\begin{exercise}[4]
Study counting-process formulations of survival analysis.
\end{exercise}

\begin{exercise}[4]
Investigate recurrent-event models.
\end{exercise}

\begin{exercise}[4]
Study shared frailty models.
\end{exercise}

\begin{exercise}[4]
Investigate inverse probability of censoring weighting.
\end{exercise}

\begin{exercise}[4]
Study interval-censored survival likelihoods.
\end{exercise}

\begin{exercise}[4]
Investigate survival calibration and time-dependent prediction metrics.
\end{exercise}

%%%%%%%%%%%%%%%%%%%%%%%%%%%%%%%%%%%%%%%%%%%%%%%%%%%%%%%%%%%%
\subsection{Conceptual Summary}
\label{subsec:survival-summary}
%%%%%%%%%%%%%%%%%%%%%%%%%%%%%%%%%%%%%%%%%%%%%%%%%%%%%%%%%%%%

Survival analysis studies the event time

\[
\boxed{
T\geq0.
}
\]

The survival function is

\[
\boxed{
S(t)
=
P(T>t).
}
\]

The hazard is

\[
\boxed{
h(t)
=
\frac{
f(t)
}{
S(t)
}.
}
\]

The cumulative hazard is

\[
\boxed{
H(t)
=
\int_0^t
h(u)\,du
=
-\log S(t).
}
\]

Therefore,

\[
\boxed{
S(t)
=
e^{-H(t)}.
}
\]

Right-censored observations are represented by

\[
\boxed{
Y_i
=
\min(T_i,C_i),
}
\]

with event indicator

\[
\boxed{
\delta_i
=
\mathbf{1}_{\{T_i\leq C_i\}}.
}
\]

The Kaplan--Meier estimator is

\[
\boxed{
\widehat{S}(t)
=
\prod_{t_j\leq t}
\left(
1-
\frac{
d_j
}{
n_j
}
\right).
}
\]

The Nelson--Aalen estimator is

\[
\boxed{
\widehat{H}(t)
=
\sum_{t_j\leq t}
\frac{
d_j
}{
n_j
}.
}
\]

The Cox proportional-hazards model is

\[
\boxed{
h
(
t\mid\mathbf{x}
)
=
h_0(t)
e^{\mathbf{x}^T\boldsymbol{\beta}}.
}
\]

A one-unit predictor change corresponds to hazard ratio

\[
\boxed{
e^{\beta_j}.
}
\]

Accelerated failure-time models instead use

\[
\boxed{
\log T
=
\mathbf{x}^T
\boldsymbol{\beta}
+
\varepsilon.
}
\]

Survival analysis therefore combines

\[
\boxed{
\text{probability}
+
\text{censoring}
+
\text{risk sets}
+
\text{hazards}
+
\text{regression}.
}
\]

%%%%%%%%%%%%%%%%%%%%%%%%%%%%%%%%%%%%%%%%%%%%%%%%%%%%%%%%%%%%
\subsection{Section Connections}
\label{subsec:survival-connections}
%%%%%%%%%%%%%%%%%%%%%%%%%%%%%%%%%%%%%%%%%%%%%%%%%%%%%%%%%%%%

\chapterconnection{
Survival Analysis extends the probability-distribution, estimation, and
regression methods developed earlier in Chapter~\ref{chap:statistics} to
time-to-event outcomes with censoring and truncation. The Kaplan--Meier
and Nelson--Aalen estimators provide nonparametric estimates of survival
and cumulative hazard, while the Cox proportional-hazards model provides
a semiparametric regression framework. Counting-process formulations
connect survival analysis with stochastic processes, and frailty models
connect it with hierarchical and random-effects modeling. The next
section, Statistical Learning, develops predictive modeling, model
selection, regularization, classification, trees, ensembles, kernels,
and modern high-dimensional statistical methods.
}

%%%%%%%%%%%%%%%%%%%%%%%%%%%%%%%%%%%%%%%%%%%%%%%%%%%%%%%%%%%%
% SECTION SOURCE TRACKING
%%%%%%%%%%%%%%%%%%%%%%%%%%%%%%%%%%%%%%%%%%%%%%%%%%%%%%%%%%%%

%------------------------------------------------------------
% 8.17 Survival Analysis
%------------------------------------------------------------
%
% Source basis:
%   - Standard survival-analysis theory.
%   - Standard event-history and reliability methods.
%   - Standard proportional-hazards and parametric survival models.
%
% Used for:
%   - event-time distributions;
%   - survival functions;
%   - density functions;
%   - hazard functions;
%   - cumulative hazards;
%   - survival-hazard relationships;
%   - mean and median survival;
%   - right, left, and interval censoring;
%   - truncation and delayed entry;
%   - censored likelihoods;
%   - informative and noninformative censoring;
%   - risk sets;
%   - Kaplan--Meier estimation;
%   - Greenwood variance;
%   - Nelson--Aalen estimation;
%   - log-rank tests;
%   - exponential survival models;
%   - Weibull survival models;
%   - log-normal and log-logistic models;
%   - proportional hazards;
%   - Cox regression;
%   - hazard-ratio interpretation;
%   - partial likelihood;
%   - tied events;
%   - baseline survival estimation;
%   - nonproportional hazards;
%   - time-dependent covariates;
%   - stratified Cox models;
%   - accelerated failure-time models;
%   - survival-model diagnostics;
%   - competing risks;
%   - cumulative incidence;
%   - recurrent events;
%   - counting processes;
%   - frailty models;
%   - discrete-time survival;
%   - reliability;
%   - restricted mean survival time;
%   - landmark analysis;
%   - immortal-time bias;
%   - interval-censored data;
%   - survival prediction;
%   - concordance and Brier-score previews;
%   - applications and exercises.
%
% Notes:
%   - Statistical Learning is developed next in Section 8.18.
%   - Stochastic-process theory from Chapter 7 provides a deeper
%     foundation for counting-process survival theory.
%   - Computational Statistics later develops numerical fitting,
%     simulation, bootstrap, and predictive validation methods.
%
%%%%%%%%%%%%%%%%%%%%%%%%%%%%%%%%%%%%%%%%%%%%%%%%%%%%%%%%%%%%